\section{Introduction}\label{sec:introduction}

Cross-national survey programs fix their items long before they see how those items behave. A question on trust in institutions or on willingness to take risks must be translated, pretested and fielded in many countries before a researcher learns whether its country-level differences follow the construct it was meant to capture. Language models offer a cheap first reading of such items \citep{argyle2023,horton2023,buskirk2026}, but whether a model deserves that role is a question about a chain of links and not about a single accuracy figure. We evaluate one instrument along the chain from an anchor, a declared country-level benchmark, through the instrument, a set of adapters trained on the signs of that benchmark, to respondents, the people whose survey answers the instrument is meant to inform. We report where along the chain the instrument holds and where it fails.

That role has drawn a literature that is mostly cautionary. Synthetic responses are sensitive to prompt wording, compress the variance of human answers, and can misportray and flatten identity groups \citep{bisbee2024,boelaert2025,wang2025}, and they depend on the order and labelling of answer options \citep{dominguezolmedo2024}. The alignment procedures that make models usable change how faithfully they simulate population attitudes \citep{lyman2025}. Methodological work in this journal concludes that validation against human data remains essential, and proposes designs that treat model output as potentially informative observations beside a human gold standard \citep{kozlowski2025,broska2025}. Most of the empirical work asks whether a model reproduces answers that people already gave, whether the model is prompted \citep{santurkar2023,durmus2024} or fine-tuned on survey response distributions \citep{cao2025,suh2025}. In the fine-tuned case the human survey sits inside the training signal. Language-model pretesting, for its part, has been studied as flaw detection in question wording \citep{sturgis2026}. Whether an instrument can say which new items will separate countries in an intended direction is a different test, and it needs a decision rule fixed before the human criterion is read.

The three links of the chain are as follows. The anchor is the Global Preferences Survey \citep[GPS;][]{falk2018}, whose country-level scores on six preference dimensions supply the signs that define what the instrument is meant to carry. The instrument is a set of parameter adapters trained on synthetic response pairs that those signs label. The respondents are the people who answered Wave~7 of the World Values Survey (WVS). A favorable result at one link does not carry to the next. This is the measurement lens that has been proposed for the evaluation of generative systems in general \citep{wallach2025}, and we apply it to one instrument under a preregistered protocol. Prompted models have already been compared with the GPS across countries and found to be weak proxies for it \citep{chen2026}. We ask the question of an instrument trained on the signs of that benchmark.

The instrument, which we call synthetic cultural agents, is built from signs alone. Each country's sign profile on the six GPS dimensions deterministically labels a shared bank of paired synthetic responses, and Direct Preference Optimization \citep{rafailov2023}, read as a regularized Bradley--Terry pairwise choice model, fits one adapter per profile. The adapters are scored on survey items with prompts that never name a country. What is specific to the construction is its training signal, which consists of six declared signs and no human answers, together with the unnamed group at evaluation. The construction is inspectable: the researcher controls which aggregate moments enter, the sign rule that turns them into labels, and whether the group is named. We built it in the development stage of this study, on a purposive panel of \dc{16}{dev_countries} countries, and evaluate it here on every country covered by both the GPS and the WVS Wave~7, with a second set of adapters for all \dc{64}{cells} possible sign profiles on a newly generated scenario bank.

The evaluation has four parts. The first asks whether training installs the synthetic labels on pairs the adapters never saw, whether the adapters' order of countries on the anchor survives a change of scenario bank, and whether a prompt that names the country or one that states the signs does as well. The second is a preregistered screen of \dc{223}{pool_items} WVS Wave~7 items that no adapter was built for. A probe that reads only the adapters' scores assigns each item a dimension and a direction, a fixed rule selects the items with the strongest assignments, and only then are the selected items compared with the human country differences. That criterion is retrospective, because the survey existed before the probe, and we say so wherever the screen is reported. The third part is a factorial over the \dc{64}{cells} sign profiles, which asks whether the six declared dimensions behave as separate inputs. The fourth consists of descriptive follow-up analyses that set the adapters beside the respondents' composite scores and answer distributions. The second-bank check, the screen and the factorial follow the preregistration. The follow-up analyses were added after the registered results had been read, and they earn no claim.

At the first link the instrument holds in a limited sense. Training installs the synthetic labels on scenarios the adapters never saw. On trust the adapters keep the GPS order of countries when the scenario bank is replaced, and a prompt that names the country does not. A prompt that states the same six signs keeps the trust order as well, and the registered contrast between adapters and that prompt is unresolved. The trust order therefore belongs to the declared signs, and the evidence cannot credit the trained weights over a prompt.

The remaining links are weaker. Each declared sign moves its own surface in the intended direction and also moves the surfaces of other dimensions, often by more, so the dimensions are entangled; a behavioral decomposition shows trust spared where the other five signs add little to its surface, and shows them carrying most of the variation on the others. Entanglement across cultural dimensions has been reported for activation steering \citep{dangmasud2026}, and the factorial here measures it for trained adapters and for a prompt. The retrospective screen of held-out items gives a weak signal: the registered null is rejected, yet the bootstrap interval reaches zero, and the registered contrasts do not distinguish the adapters from careful readers of the questionnaire on the few items they share. No dimension passes both the second-bank check and the screen. Descriptively, the adapters' order of countries does not detectably agree with the respondents' composite on any surface, and the agreement between the GPS and those respondents is itself loose on these items.

Two questions that a single agreement figure would merge separate here. One asks whether the instrument keeps the anchor's order of countries. The other asks whether it reproduces the answer distributions of respondents. On the first, the adapters keep the order on trust and a country-named prompt does not. On the second, the adapters sit farther from the survey's answer distributions than the base model does, and the country-named prompt sits closer. We report this as a descriptive contrast between two questions. The separation of group-level order from distributional fidelity is known \citep{dasilva2026}, and a country name has been found to recover a country ranking on other items in the European Social Survey \citep{wangc2026}. The crossing reported here concerns one anchor, one instrument family and the trust order of one benchmark. It was not preregistered, and it does not show that any arm matches respondents.

Our contributions are three.
\begin{enumerate}
\item \textbf{A chain-wise evaluation of one instrument.} We place the instrument between a declared anchor and a human survey and report, link by link, what a favorable or an unfavorable result permits, including the contrast between keeping the anchor's order of countries and approaching respondents' answer distributions.
\item \textbf{A preregistered retrospective screening protocol.} We specify, register and run a screen of new survey items with a probe fixed before any human value was read, a permutation null that absorbs the probe's own selection rate, a human ceiling from two survey waves, and two comparators that read only the questionnaire.
\item \textbf{A factorial over sign profiles.} We vary all six declared signs jointly to measure how far each dimension moves the surfaces of the others in trained adapters and under a prompt that states the same signs, and we add a descriptive decomposition of where trust is spared.
\end{enumerate}

What the paper shows is bounded by its design. For one anchor, one instrument family and one base model, it shows that the declared signs carry the trust order, that the instrument keeps the direction of each sign and not the separation of dimensions, that a preregistered screen finds a weak retrospective signal, and that keeping an anchor's order of countries and approaching respondents' answer distributions are different questions with different answers. It does not show that trained weights improve on a prompt that states the same signs. It does not show agreement with respondents beyond the weak retrospective signal. It does not show that screening would change which items a researcher fields, because the WVS Wave~7 file existed before the probe, the item pool and the adapters, and pre-piloting is a prospective use that remains untested. The instrument is built from signs alone, so it carries no information about magnitudes or about variation among people within a country, and its readings concern differences among countries. Within these bounds it is an inspectable instrument that a survey designer can consult beside human pretesting. It is a complement to pretesting and fielding, and the evidence gives no reason to use it in their place. Section~\ref{sec:background} places the instrument among current uses of synthetic respondents and in the measurement literature on cross-national comparability. Section~\ref{sec:instrument} describes the construction, Section~\ref{sec:design} the evaluation design and its registration, Section~\ref{sec:results} the results, and Section~\ref{sec:discussion} what they permit a researcher to do.
